\section{Existence and Uniqueness}

\subsection{Contraction Mappings and Fixed Point Theorem}

Let $(\mathcal{X}, \|\cdot\|)$ be a non-empty Banach space and $S \subset \mathcal{X}$ a
closed subset. Let $P : \mathcal{X} \to \mathcal{X}$ be a mapping that

\begin{enumerate}
    \item It is a contraction mapping on $S$, i.e., there exists $\rho \in [0, 1)$ such that
          \noindent\begin{equation*}
              \|P x - P y\| \leq \rho \|x - y\|, \quad \forall x, y \in S.
          \end{equation*}
    \item $P(S) \subset S$, i.e., $P x \in S$, for all $x \in S$.
\end{enumerate}

Then the map $P$ admits one and only one fixed point $x^*$ in $S$, meaning that
\noindent\begin{equation*}
    P x^* = x^*
\end{equation*}


\newpar{}
\ptitle{Remarks}

\begin{itemize}
    \item Banach spaces $\mathcal{X}$ are
          \begin{enumerate}
              \item complete (every Cauchy series converges to $x\in \mathcal{X}$)
              \item have a norm
          \end{enumerate}
    \item Using $P$, one can start with an arbitrary $x_0\in S$ and converges towards $x^*$.
    \item Defining $S$ so that $P$ maps $S \to S$, i.e.\ that $Px \in S$ always, is often difficult.
\end{itemize}

\subsection{Local Existence and Uniqueness}

Let $f (t, x)$ be a piecewise continuous function in $t$ and satisfy the Lipschitz condition
\noindent\begin{equation*}
    \|f (t, x)- f (t, y)\| \leq L \|x- y\|
\end{equation*}
for all $x, y \in B (x_0, r) = \{x \in \mathbb{R}^n | \|x- x_0\| \leq r\}$, $\forall t \in [t_0, t_1]$. Then there exists some $\delta > 0$ such that the state equation
\noindent\begin{equation*}
    \dot{x} = f (t, x), \quad x (t_0) = x_0
\end{equation*}
has a unique solution over $[t_0, t_0 + \delta]$.

\newpar{}
\ptitle{Autonomous Case}

\noindent\begin{equation*}
    \|f (x)- f (y)\| \leq L \|x- y\|
\end{equation*}
In simple words: If the $\Delta f$ can be bounded by a linear function within the domain of interest $B$, the function is locally Lipschitz.

\begin{examplesection}[Proof]

    If $x (t)$ is a solution of
    \noindent\begin{equation*}
        \dot{x} = f (t, x), \quad x (t_0) = x_0,
    \end{equation*}
    then by integration,
    \noindent\begin{equation*}
        x (t) = x_0 + \int_{t_0}^{t} f (s, x (s)) ds.
    \end{equation*}
    Existence and uniqueness of the solution of the differential equation is equivalent to the study of the integral equation above.

    Define the mapping (applying $P$ yields a function of time):
    \noindent\begin{equation*}
        (Px)(t) := x(t_0) + \int_{t_0}^{t} f (s, x (s)) ds, \quad t \in [t_0, t_1],
    \end{equation*}
    which maps a continuous function $x : [t_0, t_1] \to \mathbb{R}^n$ into another continuous function $\mathbb{R}^n$.

    \newpar{}
    \ptitle{Banach Space and Closed Set}

    If we can find an $x$, s.t.\ $x=(Px)$, i.e.\ $x$ is a fixed point of the above mapping, then we proved existence and uniqueness of the solution by the contraction mapping theorem (CMT). To use the CMT we must define appriopriate
    \begin{enumerate}
        \item Banach space $\mathcal{X}$
        \item closed subset $S \subset \mathcal{X}$
    \end{enumerate}
    such that $P$ maps $S$ into $S$ and is a contraction over $S$.

    We get
    \begin{enumerate}
        \item the Banach space of cont.\ functions $x:[t_0, t_0+\delta] \rightarrow \mathbb{R}^n$
              \noindent\begin{align*}
                  \mathcal{X} & = C [t_0, t_0 + \delta]                    \\
                  \|x\|_C     & = \max_{t \in [t_0, t_0+\delta]} \|x (t)\|
              \end{align*}
              where $\|x\|_C$ is the associated norm.
        \item and the closed set
              \noindent\begin{equation*}
                  S = \{x \in \mathcal{X} | \|x- x_0\|_C \leq r\}
              \end{equation*}
              with $r$ the radius of $\mathcal{B}(x_0, r)$.
    \end{enumerate}

    Note that
    \begin{itemize}
        \item $\delta$ is a positive constant to be chosen.
        \item we restrict $\delta \leq t_1-t_0$, i.e.\ $[t_0, t_0+\delta]\subset [t_0, t_1]$.
        \item $\|x\|$ is a norm on, $\mathcal{B}$ is a ball in $\mathbb{R}^n$.
        \item $\|x\|_C$ is a norm on, $S$ a ball in $\mathcal{X}$.
    \end{itemize}

    \newpar{}
    \ptitle{Showing that $P$ Maps $S$ Into $S$}

    Our definition of $P$ implies $P:\mathcal{X}\rightarrow \mathcal{X}$. Now show $S \rightarrow S$, i.e.\ show that (staying in the $S$ ball)
    \noindent\begin{equation*}
        \left\|Px-x_0\right\|_C=\max_{ t{\in}[t_0,t_0{+}\delta]}\left\|\left(Px\right)\left(t\right)-x_0\right\|\leq r
    \end{equation*}
    We use the properties
    \begin{itemize}
        \item Lipschitz
        \item $h=\max_{t\in[t_0,t_1]}\|f(t,x_0)\|$ (as $f$ is piecewise cont)
        \item $\|x\left(t\right)-x_0\|\leq r,\forall t\in[t_0,t_0+\delta], \quad \forall x \in S$
    \end{itemize}
    We can now show
    \noindent\begin{align*}
        \left\|(Px)(t)-x_0\right\| & = \left\|\int_{t_0}^t f(s,x(s))ds\right\|                                                    \\
                                   & = \left\|\int_{t_0}^t\left[f(s,x(s))-f(s,x_0)+f(s,x_0)\right]ds\right\|                      \\
                                   & \leq                        \int_{t_0}^t\left[L\left\|x\left(s\right)-x_0\right\|+h\right]ds \\
                                   & \leq                        \int_{t_0}^t\left(Lr+h\right)ds                                  \\
                                   & \leq\begin{pmatrix}
                                             t-t_0
                                         \end{pmatrix}(Lr+h)                                                                      \\
                                   & \leq\quad\delta\left(Lr+h\right)
    \end{align*}
    Hence, choosing $\delta \leq r / (Lr + h)$ ensures that $P$ maps $S$ into $S$.

    \newpar{}
    \ptitle{Contraction Mapping}

    For $x, y \in S$ we must show
    \noindent\begin{equation*}
        \| (Px)(t)- (Py)(t) \|\leq \rho \| x-y\|_C, \; \rho < 1
    \end{equation*}
    We use the same properties as before and get:
    \noindent\begin{align*}
        \| (Px)(t)- (Py)(t) \| & = \left\|\int_{t_0}^t\left[f(s,x(s))-f(s,y(s))\right]ds\right\| \\
                               & \vdots                                                          \\
                               & \leq \int_{t_0}^t L\left\|x-y\right\|_C ds                      \\
                               & = (t-t_0) L \|x-y\|_C                                           \\
                               & \leq \delta L \|x-y\|_C
    \end{align*}
    Thus,
    \noindent\begin{equation*}
        \|Px- Py\|_C \leq L\delta \|x- y\|_C \leq \rho \|x- y\|_C,
    \end{equation*}
    for $\delta < \rho/L$ with $\rho < 1$. Choosing $\delta \leq \rho/L$ ensures that $P$ is a contraction mapping over $S$.

    \newpar{}
    \ptitle{Conclusion on $S$}

    By the contraction mapping theorem, if $\delta$ satisfies:
    \noindent\begin{equation*}
        \delta \leq \min \left\{ t_1 - t_0, \frac{r}{Lr + h}, \frac{\rho}{L} \right\}, \quad \text{for } \rho < 1,
    \end{equation*}
    then the integral equation has a unique solution in $S$.

    \newpar{}
    \ptitle{Conclusion on $\mathcal{X}$}

    The aforementioned holds for functions in the ball $\mathcal{B}(x_0, r)$.

    Note that since $x(t_0) = x_0$ is inside the ball $B(x_0, r)$, any continuous solution $x(t)$ must lie inside $B(x0, r)$ for some interval of time. If we choose this time small enough, we know that any solution in $X$ lies in $S$.

    Explicitly: let $t_0 +\mu$ be the first time $x(t)$ intersects the boundary of $\mathcal{B}$ before leaving the ball. We know that (using previous properties),
    \noindent\begin{equation*}
        \left\|x\left(t\right)-x_0\right\| \leq\quad\int_{t_0}^t\left(Lr+h\right)ds
    \end{equation*}
    This can be evaluated for $t=t_0 + \mu$ yielding
    \noindent\begin{equation*}
        r=\left\|x\left(t_0+\mu\right)-x_0\right\|\leq\left(Lr+h\right)\mu
    \end{equation*}
    Therefore we know that we need at least $\mu\geq\frac{r}{Lr+h}$ units of time to leave $\mathcal{B}$ but from before we know that we choose $\delta$ even smaller than the RHS.

    Hence, any solution in $X$ lies in $S$, ensuring uniqueness in $X$ as well
\end{examplesection}


\subsubsection{Properties}
\begin{enumerate}
    \item The Lipschitz condition is stronger than continuity.
    \item Some continuous functions are not Lipschitz, e.g.,
          \noindent\begin{equation*}
              f (x) = x^{1/3}
          \end{equation*}
          which can be solved by separation of variables.
    \item The norm for Lipschitz continuity is not relevant.
    \item Locally bounded $\frac{\partial f}{\partial x}$ implies local Lipschitz continuity. Hence, check the \textbf{Jacobian} over the domain of interest first for a sufficient condition.
    \item A function can be Lipschitz but not differentiable, e.g., $f (x) = [x^2, - \text{sat}(x_1 + x_2)]$.
    \item Lipschitz is only a \textbf{sufficient} condition! % TODO: Right?
\end{enumerate}

\subsubsection{Global Existence and Uniqueness}

Let $f (t, x)$ be piecewise continuous in $t$ over $[t_0, t_1]$ and \textit{globally} Lipschitz in $x$. Then the equation:
\noindent\begin{equation*}
    \dot{x} = f (t, x), \quad x (t_0) = x_0
\end{equation*}
has a unique solution over $[t_0, t_1]$.

\subsubsection{Temporally Unbounded Existence}

Let $f(t, x)$ be piecewise continuous in $t$ and locally Lipschitz
in $x$ for $x \in D$. Let $w$ be a compact subset of $D$. Then if
\noindent\begin{equation*}
    x(t) = (Px)(t) \in w, \forall t > t_0
\end{equation*}
i.e., $w$ is invariant, then the solution $x(t)$ exists $\forall t > t_0$.

Note that the solution exists only in $w$ for all time, not in the entire domain.

\subsection{Continuity and Differentiability}

The solution of $\dot{x} = f(t, x)$ should continuously depend on $t_0$, $x_0$, $f(x,t)$.

\newpar{}
\ptitle{Dependence on Initial Time}

Follows directly from
\noindent\begin{equation*}
    x\left(t\right)=x_0+\int_{t_0}^t f\left(s,x\left(s\right)\right)ds
\end{equation*}

\newpar{}
\ptitle{Dependence on Parameters}

Let $f (t, x, \lambda)$ be continuous in $(t, x, \lambda)$ and locally Lipschitz in $x$ (uniformly in $t$ and $\lambda$) on $[t_0, t_1] \times D \times \{\|\lambda - \lambda_0\| \leq c\}$, where $D \subset \mathbb{R}^n$ is an open connected set.

Let $y (t, \lambda_0)$ be a solution of
\noindent\begin{equation*}
    \dot{x} = f (t, x, \lambda_0) \text{ with } y (t_0, \lambda_0) = y_0 \in D.
\end{equation*}

Suppose $y (t, \lambda_0)$ is defined and belongs to $D$ for all $t \in [t_0, t_1]$. Then, given $\epsilon > 0$, there exists $\delta > 0$ such that if

\noindent\begin{equation*}
    \|z_0 - y_0\| < \delta, \quad \|\lambda - \lambda_0\| < \delta,
\end{equation*}
then there is a unique solution $z (t, \lambda)$ of $\dot{x} = f (t, x, \lambda)$ defined on $[t_0, t_1]$, with $z (t_0, \lambda) = z_0$, and satisfying

\noindent\begin{equation*}
    \|z (t, \lambda) - y (t, \lambda_0)\| < \epsilon, \quad \forall t \in [t_0, t_1].
\end{equation*}


\begin{examplesection}[Proof]

    By continuity of $y (t, \lambda_0)$ in $t$ and the compactness of $[t_0, t_1]$, we know that $y (t, \lambda_0)$ is bounded on $[t_0, t_1]$.

    \newpar{}
    \ptitle{Deviation From Nominal Solution}

    Define a tube $U$ around the solution $y(t, \lambda_0)$ by
    \noindent\begin{equation*}
        U=\{(t,x)\in[t_0,t_1]\times\mathbb{R}^n|\left\|x-y\left(t,\lambda_0\right)\right\|\leq\varepsilon\}
    \end{equation*}
    We can estimate
        {\small
            \noindent\begin{equation*}
                \left\|y\left(t\right)-z\left(t\right)\right\|\leq\left\|y_0-z_0\right\|\exp\left[L\left(t-t_0\right)\right]+\frac{\mu}{L}\left\{\exp\left[L\left(t-t_0\right)-1\right]\right\}
            \end{equation*}
        }
    where $L$ is the Lipschitz constant and $\mu$ is a constant related to the size of the perturbation of the parameters.
\end{examplesection}


\subsection{Differentiability and Sensitivity Equations}

Defining
\noindent\begin{equation*}
    S(t) = {\left[\frac{\partial x}{\partial \lambda}\right]}_{\lambda = \lambda_0}
\end{equation*}
(sensitivity wrt.\ change in parameters), we can write the \textbf{sensitivity equation}:
\noindent\begin{equation*}
    \dot{S} (t) = A (t, \lambda_0) S (t) + B (t, \lambda_0), \quad S (t_0) = 0,
\end{equation*}
where
\noindent\begin{align*}
    A\left(t,\lambda\right) & = \frac{\partial f}{\partial x}\left(t,x\left(t,\lambda\right),\lambda\right)      \\
    B\left(t,\lambda\right) & = \frac{\partial f}{\partial\lambda}\left(t,x\left(t,\lambda\right),\lambda\right)
\end{align*}
are continuous functions. Hence, the continuous differentiability of $f$ with respect to $x$ and $\lambda$ implies that $x(t, \lambda)$ is differentiable with respect to $\lambda$ near $\lambda_0$.

The solution $S(t)$ tells us at which points in time $x(t,\lambda)$ will be how sensitive to change in parameters.

\newpar{}
\ptitle{Approximation of Trajectory Solution}

In a neighborhood of $\lambda_0$ we can approximate the nominal solution $x(t, \lambda_0)$ as
\noindent\begin{equation*}
    x\left(t,\lambda\right)\approx x\left(t,\lambda_0\right)+S\left(t\right)\left(\lambda-\lambda_0\right)
\end{equation*}

